\chapter{Falsification Criteria and Research Roadmap}
\label{ch:roadmap}

\section{How the model can fail}
The information-spinor interpretation should be abandoned or substantially revised if any of the following occurs:
\begin{enumerate}
\item no noncircular normalized two-branch representation of $\xi$ can be constructed;
\item every candidate representation requires insertion of $\Ree s=1/2$ before the zero condition is evaluated;
\item the complement branches can be rescaled arbitrarily, making the equal-weight conclusion gauge dependent;
\item a proposed operator is not essentially self-adjoint and no natural self-adjoint extension exists;
\item the spectral determinant reproduces only the average zero density, not the fluctuation spectrum;
\item the determinant identity contains an unanalysed prefactor with zeros in the critical strip;
\item numerical agreement disappears for frozen, out-of-sample ordinates;
\item the holonomy is a removable gauge artifact rather than a gauge-invariant cycle phase.
\end{enumerate}

These are not peripheral cautions. They define the empirical and mathematical content of the programme.

\section{Stage I: exact representation search}
The first objective is an exact zeta or Xi decomposition into complementary branches. Candidate starting points are:
\begin{enumerate}
\item the theta-function integral underlying the functional equation;
\item the Riemann--Siegel integral and approximate functional equation;
\item odd/even sectors of the eta representation;
\item Mellin transforms on $L^2(\mathbb R_+)$ with unitary dilation;
\item de Branges or Hermite--Biehler decompositions.
\end{enumerate}
Each construction must be evaluated by the no-circularity criterion of Chapter \ref{ch:conditional}.

\section{Stage II: derive equal branch normalization}
Even if two branches are found, the coefficient equality must be derived. Promising mechanisms include:
\begin{enumerate}
\item unitarity of the functional-equation scattering matrix on the fixed axis;
\item an exchange symmetry with a conserved norm;
\item a CPT-like antiunitary involution;
\item half-density normalization of the Mellin transform;
\item a projective metric that fixes the square-root amplitudes uniquely.
\end{enumerate}
A manually chosen normalization would weaken the result to a coordinate convention.

\section{Stage III: establish pointwise rather than pairwise self-duality}
This is the decisive stage. Global symmetry allows off-axis pairs. The project must identify a local principle that prevents a zero from being paired with a distinct partner. Candidate principles are:
\begin{enumerate}
\item self-adjointness of the spectral generator;
\item positivity of a Weil-type quadratic form;
\item a topological index that pins zero modes to the symmetry boundary;
\item a protected chiral zero mode in a two-component operator;
\item total positivity of the Xi kernel.
\end{enumerate}

\section{Stage IV: independent review}
Before any claim stronger than a conditional theorem, the work should undergo:
\begin{enumerate}
\item line-by-line review by an analytic number theorist;
\item operator-domain review by a functional analyst;
\item independent reproduction of the code and build;
\item adversarial search for hidden use of RH-equivalent assumptions;
\item explicit comparison with known failed proof patterns.
\end{enumerate}

\section{Publication ladder}
The work naturally separates into publishable units:
\begin{description}
\item[Paper A.] Exact equivalence of Shannon balance, complementary cancellation, Bloch equator, and Berry $-1$ holonomy.
\item[Paper B.] Dirichlet eta as a parity-phase bridge and the normalized parity representation problem.
\item[Paper C.] A candidate dilation-spinor operator with full domain and self-adjointness audit.
\item[Paper D.] Kernel positivity or de Branges formulation, if a nontrivial inequality is established.
\end{description}
Paper A can be mathematically complete without claiming RH. Papers C or D would contain the genuinely difficult number-theoretic advance.

\section{Versioning and frozen claims}
Version v0.1 freezes the following claims:
\begin{enumerate}
\item the exact bridge theorem inside the two-state model;
\item the conditional critical-axis theorem;
\item the explicit statement that the zero-state representation remains unproved;
\item the numerical audit protocol;
\item the operator and kernel proof obligations.
\end{enumerate}
Future versions must preserve this ledger append-only. A later correction must add a new record rather than rewrite the historical claim state.
